% Standalone preamble for the P29 project report.
% Deliberately NOT the notes housestyle: minimal package set, article class.
% Conventions kept from the collection standard: booktabs three-rule tables,
% \cref only, $...$ / \[...\] math delimiters, \emph emphasis, no em-dashes.

\usepackage[utf8]{inputenc}
\usepackage[T1]{fontenc}
\usepackage[english]{babel}
\usepackage[a4paper, margin=2.5cm]{geometry}

\usepackage{amsmath, amssymb}
\usepackage{array}
% Sans body, coherent with the sans labels of the exported diagrams, with
% matching sans math from newtxsf.
\usepackage[sfdefault, lining, light]{FiraSans}
\usepackage{newtxsf}
\usepackage{microtype}
\usepackage[table]{xcolor}
\definecolor{unipiblue}{RGB}{20, 74, 144} % sampled from the cherubino
\usepackage{booktabs}
\definecolor{bordeaux}{RGB}{138, 26, 45}
\definecolor{rowshade}{gray}{0.965}
% Tick and cross for the yes/no columns of tab:messages and tab:nets.
\usepackage{pifont}
\newcommand{\ck}{\textcolor{unipiblue}{\ding{51}}}
\newcommand{\cx}{\textcolor{bordeaux}{\ding{55}}}
\usepackage{graphicx}
\usepackage{tikz}
\usetikzlibrary{arrows.meta, positioning, backgrounds}
\usepackage{caption}

% Headings: hierarchy by size and whitespace, not weight. Letterspaced
% uppercase stays reserved for the running head and other small furniture.
\usepackage{titlesec}
\titleformat{\section}{\LARGE\firalight\color{unipiblue}}{\thesection}{0.8em}{}
\titleformat{\subsection}{\Large\firalight\color{unipiblue}}{\thesubsection}{0.8em}{}
\titlespacing*{\section}{0pt}{1.7em}{0.9em}
\titlespacing*{\subsection}{0pt}{1.3em}{0.6em}
\titleformat{\subsubsection}{\large\firalight\color{unipiblue}}{\thesubsubsection}{0.8em}{}
\titlespacing*{\subsubsection}{0pt}{1.0em}{0.45em}

% Block paragraphs, and open leading for the sans body.
\usepackage[skip=0.55em]{parskip}
\usepackage{setspace}
\setstretch{1.15}

\usepackage{enumitem}
\setlist{itemsep=2pt, topsep=4pt, leftmargin=*}

\usepackage{tocloft}
% No "Contents" heading: the entries start directly (babel overrides
% \contentsname at begin-document, so empty it inside babel's hook).
\addto\captionsenglish{\renewcommand{\contentsname}{}}
\setlength{\cftbeforetoctitleskip}{0pt}
\setlength{\cftaftertoctitleskip}{0pt}
\renewcommand{\cftdotsep}{\cftnodots}
\renewcommand{\cftsecfont}{\LARGE\color{unipiblue}}
\renewcommand{\cftsecpagefont}{\LARGE\color{unipiblue}}
\setlength{\cftbeforesecskip}{0.8em}
\setlength{\cftsecnumwidth}{2.2em}
\renewcommand{\cftsubsecfont}{\large}
\renewcommand{\cftsubsecpagefont}{\large}
\setlength{\cftbeforesubsecskip}{0.35em}
\setlength{\cftsubsecindent}{2.2em}
\setlength{\cftsubsecnumwidth}{2.6em}
\setcounter{tocdepth}{3}

\usepackage{fancyhdr}
\setlength{\headheight}{14pt}
\pagestyle{fancy}
\fancyhf{}
\fancyhead[L]{\footnotesize\color{unipiblue}\textls[200]{P29: THESIS}}
\fancyhead[R]{\footnotesize\firalight\color{black!60} Business Process Modelling}
\fancyfoot[L]{\footnotesize\firalight\color{black!60} Francesco Secoli \,\textperiodcentered\, 502525}
\fancyfoot[R]{\footnotesize\color{unipiblue}\thepage}
\renewcommand{\headrule}{{\color{unipiblue}\hrule width\headwidth height 0.5pt}}
\renewcommand{\footrulewidth}{0pt}

% Links in the accent blue: a cross-reference the reader can see is clickable.
\usepackage{hyperref}
\hypersetup{colorlinks=true, linkcolor=unipiblue, citecolor=unipiblue,
  urlcolor=unipiblue}

% long unbreakable paths in \texttt would otherwise overflow the measure
\emergencystretch=1em
\usepackage[capitalise, noabbrev]{cleveref}
\usepackage{etoolbox}
\AtBeginEnvironment{tabular}{\footnotesize}

% Reference and caption labels in accent blue at Fira's Medium weight, which
% reads as an accent against the Light body without the mass of Bold. Use
% \firamedium and not \fontseries{medium}: under LuaTeX the package loads Fira
% through fontspec in TU, where only m and b exist and the request falls back
% to Light in silence. The label word sits inside the hyperlink so it takes the
% colour too; a bare \crefname leaves it outside, black beside a blue number.
\newcommand{\reflabel}[1]{{\firamedium #1}}
\crefformat{table}{#2\reflabel{Table~#1}#3}
\Crefformat{table}{#2\reflabel{Table~#1}#3}
\crefrangeformat{table}{#3\reflabel{Tables~#1}#4 to #5\reflabel{#2}#6}
\crefmultiformat{table}{#2\reflabel{Tables~#1}#3}%
  { and #2\reflabel{#1}#3}{, #2\reflabel{#1}#3}{, and #2\reflabel{#1}#3}
\crefformat{figure}{#2\reflabel{Figure~#1}#3}
\Crefformat{figure}{#2\reflabel{Figure~#1}#3}
\crefrangeformat{figure}{#3\reflabel{Figures~#1}#4 to #5\reflabel{#2}#6}
\crefmultiformat{figure}{#2\reflabel{Figures~#1}#3}%
  { and #2\reflabel{#1}#3}{, #2\reflabel{#1}#3}{, and #2\reflabel{#1}#3}
\crefformat{section}{#2\reflabel{Section~#1}#3}
\Crefformat{section}{#2\reflabel{Section~#1}#3}
\crefrangeformat{section}{#3\reflabel{Sections~#1}#4 to #5\reflabel{#2}#6}
\crefmultiformat{section}{#2\reflabel{Sections~#1}#3}%
  { and #2\reflabel{#1}#3}{, #2\reflabel{#1}#3}{, and #2\reflabel{#1}#3}
\crefformat{subsection}{#2\reflabel{Section~#1}#3}
\Crefformat{subsection}{#2\reflabel{Section~#1}#3}
\crefrangeformat{subsection}{#3\reflabel{Sections~#1}#4 to #5\reflabel{#2}#6}
\crefmultiformat{subsection}{#2\reflabel{Sections~#1}#3}%
  { and #2\reflabel{#1}#3}{, #2\reflabel{#1}#3}{, and #2\reflabel{#1}#3}
\DeclareCaptionFont{medseries}{\firamedium}
\captionsetup{labelfont={color=unipiblue, medseries}}

\renewcommand{\arraystretch}{1.1}

% The assignment, quoted once in Section 1: small type, indented, in a thin
% grey box.
\newsavebox{\assignmentbox}
\newenvironment{assignment}
  {\par\medskip\begin{lrbox}{\assignmentbox}%
   \begin{minipage}{\dimexpr\linewidth-1.4cm-16pt-0.4pt\relax}%
   \footnotesize\setlength{\parskip}{0.4em}%
   }
  {\end{minipage}\end{lrbox}%
   \noindent\hspace*{0.7cm}%
   \tikz\node[draw=black!30,line width=0.2pt,rounded corners=5pt,inner sep=10pt]
     {\usebox{\assignmentbox}};%
   \par\medskip}

% Let a table take more of a page, so the two in Section 1 are not both pushed
% past the prose that introduces them.
\renewcommand{\topfraction}{0.9}
\renewcommand{\bottomfraction}{0.7}
\renewcommand{\textfraction}{0.1}
\renewcommand{\floatpagefraction}{0.8}
\setcounter{topnumber}{3}
\setcounter{bottomnumber}{2}
\setcounter{totalnumber}{4}
